\documentclass[a4paper,11pt]{article}
\usepackage[a4paper,margin=2cm]{geometry}
\usepackage[T1]{fontenc}
\usepackage[utf8]{inputenc}
\usepackage[ngerman,english]{babel}
\usepackage{lmodern,microtype,xcolor,parskip,enumitem}
\usepackage[most]{tcolorbox}
\usepackage{titlesec}
\usepackage[hidelinks]{hyperref}
\definecolor{HeaderBlueDark}{RGB}{70,110,170}
\definecolor{RowBlueLight}{RGB}{235,243,255}
\definecolor{RowBlueDark}{RGB}{220,233,250}
\definecolor{SoftGray}{RGB}{90,90,90}
\setlength{\parindent}{0pt}
\setlist[itemize]{leftmargin=1.4em,itemsep=0.35em,topsep=0.35em}
\titleformat{\section}[block]{\Large\bfseries\color{HeaderBlueDark}}{}{0pt}{}
\newtcolorbox{exbox}{enhanced,breakable,colback=RowBlueLight,colframe=RowBlueDark,boxrule=0.4pt,arc=1.5mm,left=1.2mm,right=1.2mm,top=1mm,bottom=1mm,title={Beispiele \textnormal{(Examples)}},fonttitle=\bfseries,colbacktitle=RowBlueDark,coltitle=black}
\NewDocumentEnvironment{grammarentry}{mm+b}{\begin{tcolorbox}[enhanced,breakable,colback=white,colframe=HeaderBlueDark,title={#1\hfill\normalfont\footnotesize #2},colbacktitle=HeaderBlueDark,coltitle=white,fonttitle=\bfseries]#3\end{tcolorbox}}{}
\newcommand{\GermanText}[1]{\textbf{Deutsch: }#1\par}
\newcommand{\EnglishText}[1]{{\small\color{SoftGray}\textbf{English: }#1\par}}
\newcommand{\ExampleItem}[2]{\item \textbf{#1}\\{\small\color{SoftGray}#2}}
\title{Die Modalpartikel \glqq bloß\grqq\\{\large\textnormal{The modal particle bloß}}}
\date{}
\begin{document}
\maketitle
\tableofcontents
\section{Bedeutung und Gesprächsfunktion \textnormal{(Meaning and pragmatic function)}}
\begin{grammarentry}{bloß}{Modalpartikel / Abtönungspartikel}
\GermanText{Bloß verstärkt emotionale Fragen, dringende Warnungen und unerfüllte Wünsche.}
\EnglishText{Bloß intensifies emotional questions, urgent warnings and unfulfilled wishes.}
\end{grammarentry}
\section{Satzart, Stellung und Betonung \textnormal{(Sentence mood, position and stress)}}
\begin{grammarentry}{Verwendung}{Modalpartikel / Abtönungspartikel}
\GermanText{Besonders in W-Fragen, negativen Imperativen und Wunschsätzen. Je nach Ton zeigt es Sorge, Ratlosigkeit, Dringlichkeit oder Warnung.}
\EnglishText{Especially in wh-questions, negative commands and wishes. Depending on tone it expresses worry, helplessness, urgency or warning.}
\end{grammarentry}
\GermanText{Als Abtönungspartikel ist bloß unflektierbar und verlangt keinen Kasus. Sie steht normalerweise im Mittelfeld und kann das Vorfeld eines normalen Aussagesatzes nicht allein besetzen.}
\EnglishText{As a discourse particle, bloß is invariable and governs no case. It normally occurs in the middle field and cannot stand alone in the prefield of a declarative sentence.}
\section{Typische Beispiele \textnormal{(Typical examples)}}
\begin{exbox}\begin{itemize}
\ExampleItem{Was ist bloß passiert?}{What on earth happened?}
\ExampleItem{Wie soll ich das bloß schaffen?}{How on earth will I manage that?}
\ExampleItem{Mach bloß keinen Fehler!}{Whatever you do, don't make a mistake!}
\ExampleItem{Wenn er bloß zurückkäme!}{If only he came back!}
\end{itemize}\end{exbox}
\section{Andere grammatische Rollen \textnormal{(Other grammatical roles)}}
\begin{grammarentry}{Abgrenzung}{Modalpartikel / Abtönungspartikel}
\GermanText{Als Fokuspartikel bedeutet bloß nur: Ich habe bloß fünf Euro. Diese begrenzende Funktion ist nicht modal.}
\EnglishText{As a focus particle, bloß means only: I have only five euros. This limiting meaning is not modal.}
\end{grammarentry}
\begin{exbox}\begin{itemize}
\ExampleItem{Ich habe bloß zwei Minuten.}{I have only two minutes.}
\ExampleItem{Pass bloß auf!}{You'd better watch out!}
\end{itemize}\end{exbox}
\section{Lernhinweis \textnormal{(Learning tip)}}
\GermanText{Nicht die Wortform allein, sondern der Kontext entscheidet, ob bloß als Modalpartikel verwendet wird. Die englische Übersetzung kann die kommunikative Absicht statt einer wortwörtlichen Entsprechung ausdrücken.}
\EnglishText{The word form alone does not determine whether bloß is a modal particle. English translation may express its conversational intent instead of using a word-for-word equivalent.}
\section{Quellen \textnormal{(Sources)}}
\begin{itemize}
\item IDS Mannheim: \url{https://grammis.ids-mannheim.de/terminologie/2}
\item Duden: \url{https://www.duden.de/sprachwissen/sprachratgeber/Modalpartikeln}
\end{itemize}
\end{document}
